\section{Theory: From Local Support to Certified Alarms}\label{sec:theory}

The theory has one main message. Transport identifies the target residual, but local support determines whether the corresponding alarm can be certified.

\subsection{The Identifiability Boundary}

The first result explains why local overlap is not merely a numerical diagnostic. If the target cell has no source support, the target residual in that cell is not identified by source labels and target covariates.

\begin{proposition}[No local overlap, no certification]\label{prop:no-overlap}
Let \(A\subseteq\cX\) be a subgroup-bin cell with \(Q_X(A)=\pi>0\) and \(P_X(A)=0\). For any \(\Delta\in(0,1)\), there exist two covariate-shift worlds \((P_0,Q_0)\) and \((P_1,Q_1)\) such that the observable law of the labeled source sample and unlabeled target covariates is identical in the two worlds, but
\[
r_A^{Q_1}-r_A^{Q_0}=\Delta .
\]
Consequently, no audit using only those observations can uniformly certify target reliability in \(A\).
\end{proposition}

The point is simple. Target covariates can reveal that a cell exists in deployment, but if source labels never appear there, the conditional label law inside that cell is unconstrained by the audit data. A finite-sample version is immediate: when \(P_X(A)=O(1/n)\), there is constant probability that the source audit set contains no label from \(A\). Honest inference must remain wide on that event.

\subsection{The Main Certificate}

Now suppose the cell has observed support. Write
\[
Z_i=Y_i-s(X_i),
\qquad
m_i=\EE[Z_i\given X_i].
\]
For any nonnegative audit weights \(v_i\), define
\[
\widehat r_c(v)
=
\frac{\sum_i v_iA_c(X_i)Z_i}{\sum_i v_iA_c(X_i)}
\]
and
\[
\widehat{\ess}_c(v)
=
\frac{\bigl(\sum_i v_iA_c(X_i)\bigr)^2}
{\sum_i v_i^2A_c(X_i)}.
\]
The conditional target of this weighted mean is
\[
\bar r_c(v)
=
\frac{\sum_i v_iA_c(X_i)m_i}{\sum_i v_iA_c(X_i)} .
\]
This is the right finite-sample object when the source covariates and target covariates have been observed. Population certification adds one more transport term below.

\begin{assumption}[Label-independent audit weights]\label{ass:label-independent-weights}
Conditional on the source and target covariates and on any randomness used to fit the weights, the audit weights \(v_i\) are fixed and the labels \(Y_i\) are independent. This holds for oracle weights \(v_i=w(X_i)\). It also holds for learned weights \(v_i=\widehat w_i\) when the weight assigned to source point \(i\) is trained without using \(Y_i\), for example by source-side cross-fitting.
\end{assumption}

Learned weights also require a local nuisance bound. For each cell, let
\[
S_c=\sum_i w(X_i)A_c(X_i),
\qquad
\Delta_c=\sum_i|\widehat w_i-w(X_i)|A_c(X_i).
\]
Assume a conservative envelope \(\rho_c\) satisfies
\begin{equation}\label{eq:rho-envelope}
\rho_c\ge \frac{2\Delta_c}{S_c-\Delta_c},
\qquad S_c>\Delta_c .
\end{equation}
For oracle weights, set \(\rho_c=0\).

\begin{theorem}[Locally certified subgroup-bin inference]\label{thm:main-certificate}
Assume \cref{ass:label-independent-weights}. Let \(M=|\cC|\), \(Y_i\in\{0,1\}\), \(s(X_i)\in[0,1]\), and suppose all audited denominators are positive. Then, conditional on the covariates and fitted weights,
\[
\PP\left(
\forall c\in\cC:
\left|\widehat r_c(v)-\bar r_c(v)\right|
\le
\sqrt{\frac{\log(2M/\delta_1)}{2\widehat{\ess}_c(v)}}
\right)
\ge
1-\delta_1 .
\]
If \(v_i=\widehat w_i\) are learned weights satisfying \cref{eq:rho-envelope}, then on the same event,
\[
\forall c\in\cC:
\left|\widehat r_c(\widehat w)-\bar r_c(w)\right|
\le
\sqrt{\frac{\log(2M/\delta_1)}{2\widehat{\ess}_c(\widehat w)}}
\;+\;\rho_c .
\]
Finally, suppose \(0\le w(X)\le W\). Let
\[
\epsilon_n(\delta_2)=2W\sqrt{\frac{\log(4M/\delta_2)}{2n}},
\qquad
\gamma_c(\delta_2)=
\frac{2\epsilon_n(\delta_2)}{\mu_c^Q-\epsilon_n(\delta_2)}
\]
for cells with \(\mu_c^Q>\epsilon_n(\delta_2)\). Then, with probability at least \(1-\delta_1-\delta_2\),
\[
\forall c\in\cC:
\left|\widehat r_c(\widehat w)-r_c^Q\right|
\le
\underbrace{
\sqrt{\frac{\log(2M/\delta_1)}{2\widehat{\ess}_c(\widehat w)}}
}_{\textnormal{local label noise}}
+
\underbrace{\rho_c}_{\textnormal{weight nuisance}}
+
\underbrace{\gamma_c(\delta_2)}_{\textnormal{source design}}.
\]
The oracle-weight case is recovered by taking \(\widehat w=w\) and \(\rho_c=0\).
\end{theorem}

The theorem separates three sources of uncertainty. The first is the unavoidable label noise in the selected local cell, and its scale is exactly local ESS. The second is density-ratio error; cross-fitting makes the label-noise argument valid but does not make \(\rho_c\) small. The third is the difference between the observed finite source covariate design and the population target residual \(r_c^Q\). In many audits the conditional guarantee is the most transparent report; population claims require the additional \(\gamma_c\) term or a sharper empirical-process analysis.

The proof has the same structure. Conditional on the covariates and fitted weights, the estimator is a weighted average of independent Bernoulli residuals. Hoeffding's inequality gives a radius controlled by the sum of squared normalized weights, which is \(1/\widehat{\ess}_c\). The learned-weight part is a deterministic ratio perturbation: if the estimated weights differ from the oracle weights by \(\Delta_c\) inside the cell, the conditional target can move by at most \(2\Delta_c/(S_c-\Delta_c)\). The population term applies the same ratio perturbation to the empirical source design.

The nuisance term is best interpreted as a local robustness penalty, not as another sampling variance. It asks: how much could the cellwise residual target move if the learned density ratio is replaced by a plausible oracle ratio? When a model-based log-ratio error bound is available, the theorem gives a direct conversion. If
\[
\sup_{i:A_c(X_i)=1}
\left|\log \widehat w_i-\log w(X_i)\right|
\le \kappa_c ,
\qquad \kappa_c<\log 2,
\]
then
\[
\rho_c
\le
\frac{2(e^{\kappa_c}-1)}{2-e^{\kappa_c}} .
\]
In most deployments, such a bound is not available from first principles. The practical substitute is a sensitivity envelope: fit a collection \(\mathcal W\) of cross-fitted density-ratio models using different regularization levels, clipping thresholds, feature sets, and resampling perturbations, and compute
\[
\widehat\rho_c^{\mathrm{sens}}
=
\max_{\widetilde w\in\mathcal W}
\frac{
2\sum_i|\widehat w_i-\widetilde w_i|A_c(X_i)
}{
\sum_i\widetilde w_iA_c(X_i)
-
\sum_i|\widehat w_i-\widetilde w_i|A_c(X_i)
},
\]
over candidates with positive denominator. This is a valid \(\rho_c\) only if the candidate envelope contains the oracle ratio with high probability; otherwise it is a transparent sensitivity analysis. If \(\widehat\rho_c^{\mathrm{sens}}\) is comparable to the observed gap, the audit should report insufficient evidence rather than a target-certified alarm.

This also gives a useful scale calculation. Ignoring nuisance and population terms, to certify a gap \(g=|\widehat r_c|\) above tolerance \(\tau\), one needs roughly
\[
\widehat{\ess}_c
\gtrsim
\frac{\log(2M/\delta)}{2(g-\tau)^2}.
\]
For example, with \(M=150\), \(\delta=0.10\), and a margin \(g-\tau=0.20\), the required local ESS is about \(100\). A cell supported by ten effective labels cannot certify such a margin, no matter how large the global ESS is.

\begin{corollary}[Certified alarm rule]\label{cor:certified-alarm}
Let
\[
B_c
=
\sqrt{\frac{\log(2M/\delta_1)}{2\widehat{\ess}_c(\widehat w)}}
\;+\;\rho_c+\gamma_c(\delta_2).
\]
If \(|r_c^Q|\le \tau\) for all \(c\in\cC\), then the rule
\[
\exists c\in\cC:
|\widehat r_c(\widehat w)|-B_c>\tau
\]
has false-alarm probability at most \(\delta_1+\delta_2\), whenever the nuisance envelope is valid.
\end{corollary}

This is the paper's practical certificate. A point estimate becomes an actionable alarm only after subtracting a simultaneous local uncertainty radius. If the point estimate exceeds \(\tau\) but the certified excess is zero, the audit has not shown that the subgroup is safe; it has shown that the current evidence is insufficient.

\subsection{Why Global ESS Cannot Replace Local ESS}

The global ESS can be large while the selected cell is effectively supported by one weighted label.

\begin{proposition}[Large global ESS with degenerate local ESS]\label{prop:global-misleading}
For any target cell mass \(\pi\in(0,1)\) and integer \(N\ge2\), there is a weighted sample such that the cell has weighted mass \(\pi\) and local ESS equal to \(1\), while the global ESS equals
\[
\frac{1}{\pi^2+(1-\pi)^2/(N-1)}.
\]
\end{proposition}

For \(\pi=0.05\), this quantity approaches \(400\) as \(N\) grows. Thus a global ESS near \(400\) does not rule out a worst-cell alarm supported by a single effective label.

\subsection{Search, Stronger Estimators, and the Same Bottleneck}

\Cref{thm:main-certificate} remains valid after selecting the empirically worst cell from a pre-specified finite family \(\cC\), because the bound holds simultaneously over that family. If analysts search over a richer class after seeing labels, they must either certify over the full searched family or reserve fresh labels for the final selected cell. Sample splitting is simple: use one part of the source audit labels to discover a candidate cell, and an independent part to apply the one-cell version of the certificate.

Doubly robust estimators can reduce nuisance bias when an outcome regression is accurate, but they do not remove the local-overlap frontier. Their influence functions still contain a cell-local weighted residual term. When a target subgroup-bin cell lacks source labels, no augmentation can identify its target conditional label law without additional assumptions.
